\section{NVMe 队列如何把块请求交给 SSD}

NVMe 控制器通过 PCIe 接入系统。主机先在 DRAM 中建立 submission queue（SQ）与 completion queue（CQ），把它们的物理地址和深度写入控制器寄存器。提交一次读请求时，软件先填好 SQ 项，其中包含 opcode、namespace、逻辑块范围和目标缓冲区描述；随后执行必要的内存屏障，再写 MMIO doorbell 更新队尾。

控制器看到 doorbell 后 DMA 读取命令和数据指针，访问 NAND/FTL，最后把完成项 DMA 写入 CQ，再更新完成相位并选择轮询或 MSI-X 通知。doorbell 写返回只说明寄存器事务完成，不代表数据已经到达内存。

\begin{lstlisting}
host memory: write SQ entry -> memory barrier -> ring doorbell
controller:  DMA command -> execute media work -> DMA data
completion:  write CQ entry -> interrupt/poll -> host consumes CQ
\end{lstlisting}

PRP 或 SGL 描述数据缓冲区的物理页列表。它们必须在命令有效期内保持固定并允许设备访问；IOMMU 映射、页锁定和完成后的解映射共同划定缓冲区所有权。队列深度提高并行度，但超过介质与控制器能并行处理的程度后，只会增加排队和尾延迟。

\section{错误、复位与持久化边界}

完成项中的状态码、命令标识和 SQ 标识决定主机能否把结果归还给正确请求。超时后不能立刻重用命令缓冲，因为旧控制器事务可能仍在路上；驱动通常先停止新提交、屏蔽通知、确认或复位队列，再按协议处理未完成请求。

普通写完成可能只表示数据进入控制器易失缓存。需要持久化时，文件系统或数据库还要使用 FUA、FLUSH 等语义，并让控制器的断电保护或介质提交边界参与保证。性能数字必须说明测的是提交、完成还是持久化完成。

\section*{章末练习}

\begin{longtable}{L{0.08\linewidth}L{0.32\linewidth}L{0.5\linewidth}}
\toprule
题号 & 类型 & 题目 \\
\midrule
1 & Flash 约束 & 为什么 NAND 不能像 DRAM 一样原地覆盖任意字节？说明 page、block 与擦除的关系。 \\
2 & FTL & 顺序逻辑写经过 FTL 后一定形成顺序物理写吗？列出映射、垃圾回收和磨损均衡的影响。 \\
3 & NVMe trace & 从主机填写 SQ 项开始，按顺序追踪一次 NVMe 读到 CQ 完成，并标出 CPU 与控制器的 DMA 所有权交接。 \\
4 & 队列 & SQ/CQ 为什么需要命令标识和 phase？队列环回时只比较数组下标有什么问题？ \\
5 & 超时恢复 & 命令超时后为什么不能立即释放数据页和命令标识？给出安全复位次序。 \\
6 & 持久化 & 普通写完成、FUA 写完成与 FLUSH 完成分别能支持什么推断？ \\
\bottomrule
\end{longtable}

\section*{参考解答与要点}

\begin{longtable}{L{0.08\linewidth}L{0.32\linewidth}L{0.5\linewidth}}
\toprule
题号 & 类型 & 解答要点 \\
\midrule
1 & Flash 约束 & NAND 以 page 编程、以包含多个 page 的 block 擦除；改写旧数据通常写到新 page，再由垃圾回收搬迁有效页并擦除整块。 \\
2 & FTL & 不一定。逻辑到物理映射可重定向位置，垃圾回收搬移有效页，磨损均衡还会主动迁移数据；这些动作改变物理顺序和写放大。 \\
3 & NVMe trace & 主机写 SQ 与数据描述并同步，写 doorbell；控制器 DMA 取命令，访问 FTL/NAND，把数据 DMA 到主机页，写 CQ 与 phase；主机在中断/轮询后同步并回收。 \\
4 & 队列 & 命令标识把乱序完成归属到请求，phase 区分环回后的新旧 CQ 项；仅看下标会把上一圈残留内容误作新完成。 \\
5 & 超时恢复 & 控制器可能仍在读写旧页并返回迟到完成。先停新提交、屏蔽通知、取消或复位队列并确认 DMA 停止，再按 generation 处理旧请求，最后解映射和复用资源。 \\
6 & 持久化 & 普通完成可能只到易失缓存；FUA 要求该写跨过设备定义的非易失边界；FLUSH 要求此前相关写到达该边界。它们仍需上层正确安排数据与元数据顺序。 \\
\bottomrule
\end{longtable}
